\chapter{Conclusão}
\label{cap:conclusão}

Este trabalho apresentou o desenvolvimento e a avaliação de um sistema de correção automatizada de provas discursivas baseado em múltiplos agentes de Inteligência Artificial, com suporte à Recuperação Aumentada por Geração (RAG) e mecanismo de arbitragem condicional por divergência. O objetivo principal foi implementar um protótipo funcional capaz de executar o fluxo completo, desde a configuração da prova até a disponibilização dos resultados para revisão docente, produzindo artefatos estruturados e revisáveis: nota final, notas por critério e indicação de divergência entre corretores.

Os resultados obtidos no estudo de caso controlado, com dados pareados entre as condições experimentais, indicam que:

\begin{enumerate}
    \item \textbf{QA1 --- Execução \textit{end-to-end}:} o sistema completou todas as etapas do fluxo sem falhas nas duas condições testadas, processando e persistindo 100\% das respostas previstas (24/24), com completude estrutural de 100\%. Cada correção produziu nota final, notas individuais de C1 e C2, valor de divergência e método de consenso, atendendo plenamente ao critério de sucesso definido na metodologia.

    \item \textbf{QA2 --- Utilidade do RAG:} a comparação pareada entre Condição A (com RAG) e Condição B (sem RAG), utilizando o mesmo conjunto de respostas em ambas as condições, revelou que o RAG aumenta a especificidade avaliativa, com efeito mais expressivo no nível intermediário ($\Delta = +0{,}84$ na média), especialmente em Q1 ($\Delta = +1{,}92$). No nível fraco, o efeito foi modesto na média ($\Delta = +0{,}12$), com casos de leve aumento de rigor (Q3: $\Delta = -0{,}55$), evidenciando que o material recuperado funciona como referencial que restringe o espaço de interpretação do corretor. Os extremos (``Excelente'' e ``Fora do tema'') não foram afetados ($\Delta = 0{,}00$).

    \item \textbf{QA3 --- Arbitragem por divergência:} em nenhum dos 24 casos avaliados (12 por condição) a divergência entre C1 e C2 ultrapassou o limiar de $2{,}0$ pontos. O valor máximo observado foi $1{,}22$ pontos. O árbitro não foi acionado, e o método de consenso \texttt{mean\_2} foi aplicado em todas as correções. O critério de sucesso (árbitro acionado somente quando $|$C1$-$C2$| > 2{,}0$) foi atendido integralmente: o mecanismo calculou divergências corretamente e não realizou acionamento indevido.

    \item \textbf{QA4 --- Estabilidade operacional:} nas três repetições ($R=3$) sobre a questão Q2, os níveis ``Excelente'' e ``Fora do tema'' apresentaram variação nula ($0{,}00$ ponto), demonstrando reprodutibilidade perfeita nos extremos da escala. O nível ``Intermediária'' (variação máxima de $0{,}84$) atendeu ao critério de $\leq 1{,}0$ ponto. O nível ``Fraca'' (variação máxima de $1{,}37$, DP $= 0{,}75$) ultrapassou o limiar, refletindo a ambiguidade inerente a respostas com erros conceituais parciais. A hierarquia de médias ($10{,}00 > 9{,}33 > 1{,}37 > 0{,}00$) confirma a capacidade discriminativa do sistema.
\end{enumerate}

Do ponto de vista arquitetural, a combinação de FastAPI para o \textit{backend}, LangGraph para orquestração do pipeline com corretores paralelos, LangChain para integração com LLMs, ChromaDB para armazenamento vetorial e PostgreSQL para persistência relacional mostrou-se viável e funcionalmente coerente. A comunicação entre o \textit{notebook} de experimentação e o sistema ocorreu integralmente via API REST, o que garantiu que os resultados refletissem o comportamento real do sistema em produção.

Do ponto de vista operacional, o sistema preserva a autoridade docente ao posicionar a correção automatizada como etapa de apoio: o professor tem acesso às notas individuais de cada corretor, ao valor de divergência e à nota final consensuada, podendo validar ou ajustar cada resultado antes da consolidação. O tempo de execução para as 12 respostas da Condição A foi de $155$ segundos ($\approx 12{,}9$ s/resposta), e os três \textit{runs} do QA4 foram consistentes ($55$ s cada), indicando estabilidade operacional do pipeline.

É importante reconhecer as limitações do estudo. O escopo restrito (uma disciplina, três questões, respostas sintéticas em quatro níveis bem definidos) não permite generalizações amplas sobre eficácia educacional ou desempenho em contextos reais. A ausência de gabarito humano como referência impede afirmar que as notas atribuídas pelo sistema são ``corretas'', apenas que são internamente consistentes e estruturalmente completas. Adicionalmente, o mecanismo de arbitragem não foi exercitado empiricamente neste experimento, pois nenhuma divergência ultrapassou o limiar configurado; sua validação em condições reais fica condicionada a estudos com maior volume e heterogeneidade de respostas. Os resultados devem ser interpretados como evidência de viabilidade técnica e operacional do protótipo, não como validação pedagógica conclusiva.

\section{Trabalhos Futuros}

Com base nos resultados e limitações identificados, sugere-se as seguintes direções para trabalhos futuros:

\begin{itemize}
    \item \textbf{Validação com respostas humanas reais:} aplicar o sistema em turmas reais para avaliar a concordância entre notas atribuídas pela IA e notas atribuídas por docentes, utilizando métricas como coeficiente Kappa e correlação de Pearson.

    \item \textbf{Exercitação empírica do árbitro:} ampliar o conjunto de respostas e diversificar os perfis de qualidade de modo a provocar divergências acima do limiar, permitindo validar o mecanismo de arbitragem em condições reais e avaliar o impacto da regra de consenso \textit{closest-pair} sobre a nota final.

    \item \textbf{Estudo comparativo multi-modelo:} avaliar o desempenho do pipeline com diferentes provedores de LLM (OpenAI, Anthropic, modelos locais via Ollama), comparando qualidade das correções, estabilidade operacional e custo por resposta.

    \item \textbf{Análise aprofundada do impacto do RAG:} expandir a comparação A vs.\ B com maior diversidade de questões, disciplinas e materiais didáticos, investigando em quais contextos o efeito de especificidade avaliativa proporcionado pelo RAG é mais significativo e consistente.

    \item \textbf{Interface de revisão docente completa:} expandir a interface para incluir ajuste de nota por critério, anotações do professor e relatórios pedagógicos que identifiquem lacunas de aprendizado e tendências da turma.

    \item \textbf{Escalabilidade e implantação:} avaliar o sistema em cenários com dezenas de alunos e múltiplas provas simultâneas, otimizando gerenciamento de \textit{rate limits} e custos de API.

    \item \textbf{Detecção de plágio e integridade acadêmica:} integrar mecanismos de detecção de similaridade entre respostas e de conteúdo gerado por IA, contribuindo para a integridade dos processos avaliativos.
\end{itemize}

\glsaddall
